%!TEX program = lualatex

% ============================================================================
% flat_profile_theory.tex  —  самодостаточная сборка (xelatex, дважды):
%   xelatex -halt-on-error -interaction=nonstopmode flat_profile_theory.tex
% ============================================================================
\documentclass[11pt,a4paper]{article}
\usepackage{fontspec}
\setmainfont{Noto Serif}
\setsansfont{Noto Sans}
\setmonofont{DejaVu Sans Mono}[Scale=MatchLowercase]
\usepackage[russian]{babel}
\usepackage[margin=21mm]{geometry}
\usepackage{amsmath,amssymb,amsthm,mathtools,booktabs,tabularx,array}
\usepackage{microtype,enumitem,xcolor}
\usepackage[colorlinks=true,linkcolor=blue!45!black,citecolor=blue!45!black,urlcolor=blue!45!black]{hyperref}
\setlength{\parindent}{0pt}
\setlength{\parskip}{4pt}
\setlength{\emergencystretch}{2em}
\setlist{nosep,leftmargin=*}
\setcounter{tocdepth}{1}
\newcommand{\R}{\mathbb R}
\newcommand{\Id}{\mathrm I}
\newcommand{\T}{^{\mathsf T}}
\newcommand{\F}{_{\mathrm F}}
\newcommand{\ip}[2]{\langle #1,#2\rangle_{\mathrm F}}
\newcommand{\Gr}{\mathrm{Gr}}
\newcommand{\Exp}{\operatorname{Exp}}
\newcommand{\grad}{\operatorname{grad}}
\newcommand{\qf}{\operatorname{qf}}
\newcommand{\rot}{\operatorname{rot}}
\DeclareMathOperator{\rank}{rank}
\DeclareMathOperator{\tr}{tr}
\DeclareMathOperator{\diag}{diag}
\DeclareMathOperator{\spanop}{span}
\DeclareMathOperator*{\argmin}{arg\,min}
\newtheorem{theorem}{Теорема}[section]
\newtheorem{proposition}[theorem]{Утверждение}
\newtheorem{corollary}[theorem]{Следствие}
\newtheorem{definition}[theorem]{Определение}
\theoremstyle{remark}
\newtheorem{remark}[theorem]{Замечание}
\newcolumntype{L}{>{\raggedright\arraybackslash}X}

\title{Профильно-плоская схема\\[3pt]
\large Единая теория SVD-решателя, Grassmann-спуска\\
и анизотропной адаптации для многоиндексной задачи}
\author{Исследовательская записка по материалам \texttt{SVD/},
\texttt{grassmann\_rank\_one}, \texttt{cladue.md}, \texttt{Tucker.md}}
\date{октябрь 2026 г.}

\begin{document}
\maketitle
\vspace{-8mm}
\begin{abstract}
Интуиция «разложить подпространство или изменённое многообразие на простое,
быстро вычислимое» формализуется как схема из трёх операций:
(О1) \emph{профиль} --- исключение линейных параметров, снимающее
калибровочную неоднозначность; (О2) \emph{модификация многообразия} ---
карта или однопараметрическая деформация экспоненты
 $R^\varphi_Y=\Exp_Y(V\varphi(\Sigma)A\T)$;
(О3) \emph{плоское разложение} --- картановское разложение шага на
попарно коммутирующие двумерные вращения с точным скалярным поиском
в каждом. На симметрических пространствах (евклидово пространство матриц,
 $\Gr(m,d)$, конус PSD ранга $m$) разложение точное.
SVD-решатель --- евклидов частный случай: rank-one поправка с
QR-нормировкой геометрически совпадает с rank-one геодезическим шагом,
изотропная модель $\mathcal LX=GXH$ решается замкнуто (Эккарт--Янг),
жадный выбор атома и точный шаг $t_*=R/D$ --- спектральный бюджет и
скалярный профиль схемы. Сформулированы гарантии спуска и сходимости,
посчитаны стоимости и предложены улучшения: плоский шаг ранга $q$ с совместным GN по углам, предобусловленный выбор плоскостей,
рациональная (полярная) параметризация профиля, $\kappa$-гомотопия
OPG$\to$MAVE, сжатие статистик, двухуровневая доводка.
\end{abstract}
\vspace{-3mm}
\tableofcontents
\medskip
\textbf{Статус утверждений.} Все формулы выведены; доказательства либо
приведены, либо даны ссылкой. Предложения алгоритмов помечены явно;
сквозные характеристики (wall-clock, качество восстановления) требуют
замера и не заявляются.

\section{Модель и функционалы}\label{sec:model}

Пусть $y_i=\psi(Y_*\T x_i)+\varepsilon_i$, $Y_*\T Y_*=\Id_m$, $m\ll d$;
при $\mathbb E[\nabla\psi\,\nabla\psi\T]\succ0$ пространство градиентов
совпадает с EDR-пространством. Для центра $x_j$:
 $w_{ij}=K\bigl((x_i-x_j)\T D_j(x_i-x_j)\bigr)$, $N_j=\sum_iw_{ij}$,
 $\Sigma_j=N_j^{-1}\sum_iw_{ij}z_{ij}z_{ij}\T$, $c_j=N_j^{-1}\sum_iw_{ij}z_{ij}y_i$ с $z_{ij}=x_i-\bar x_j$; для матрицы направлений $S_j$:
 $U_j=S_j\Sigma_j\in\R^{p\times d}$, $I_j=S_jc_j$, локальная оценка градиента
 $\hat\nabla_j=\Sigma_j^{-1}c_j$ (в $area{\spanop}$ невырожденности).
Локальная модель: $I_j\simeq U_jYg_j$, $g_j\in\R^m$, $Y\T Y=\Id_m$.

Четыре различные роли метрик (путать нельзя):
(1) $D_j$ --- локализация, меняет статистику;
(2) $W_j\succeq0$ --- метрика невязки;
(3) $H\succ0$ --- метрика штрафа;
(4) $\mathcal M_Y\succ0$ на касательном --- метрика шага.
Далее статистики заморожены; изменение $D_j,S_j,h$ начинает новую задачу.

Профилированный целевой функционал на $\Gr(m,d)$:
\begin{equation}\label{eq:profile}
 \Phi(Y)=\sum_jN_j\min_{g\in\R^m}\bigl\{\|I_j-U_jYg\|^2+\tau\|g\|^2\bigr\}
 +\lambda\,d_{\rm ch}^2(Y,Y_0),\qquad
 d_{\rm ch}^2=\tfrac12\|YY\T-Y_0Y_0\T\|\F^2=m-\|Y_0\T Y\|\F^2.
\end{equation}
При $\tau>0$ профиль аналитичен на компактном $\Gr$; при $\tau=0$ он лишь
полунепрерывен сверху и минимум может не достигаться.
При $\tau=0$ и без штрафа свободная полноранговая матрица и её
ортонормированный базис дают один профиль ($GL(m)$-замена коэффициентов);
штраф $\|B-P\|\F^2$ по базисам этим свойством не обладает.

\section{Профильно-плоская схема}\label{sec:scheme}

\subsection{Разделимый класс и профиль}\label{sec:master}

\begin{definition}[разделимый класс]\label{def:master}
\begin{equation}\label{eq:master}
 \Phi(\theta)=\min_c\bigl\{\|b-\mathcal A(\theta)c\|^2+r(c)\bigr\}+\Omega(\theta),
\end{equation}
где $\theta$ --- структура (матрица, подпространство, тензор, поле),
 $c$ --- линейные коэффициенты, $\mathcal A$ --- эквивариантный словарь.
\end{definition}

\begin{proposition}[корректность профиля и калибровка]\label{prop:profile}
Если $\widetilde{\mathcal A}=[\mathcal A;\mathcal R_{r(c)}]$ имеет локально
постоянный столбцовый ранг, то
 $\Phi=\|\tilde b\|^2-\tilde b\T\widetilde{\mathcal A}
(\widetilde{\mathcal A}\T\widetilde{\mathcal A})^{-1}\widetilde{\mathcal A}\T\tilde b+\Omega$,
 $\tilde b=[b;0]$. Если при $\theta\mapsto\theta Q$ словарь преобразуется как
 $\Psi_j(\theta Q)=\Psi_j(\theta)\rho(Q)$, а $r,\Omega$ инвариантны при
 $c_j\mapsto\rho(Q)^{-1}c_j$, то \eqref{eq:master} корректно определён на
классе эквивалентности $[\theta]$.
\end{proposition}
\begin{proof}
Первая формула --- исключение квадратичной формы по $c$. Вторая: замена
коэффициентов взаимно однозначна и сохраняет каждое значение минимизируемого
выражения.
\end{proof}

Частные случаи \eqref{eq:master}: fixed-$g$ SVD-солвер (калибровка тривиальна,
 $c$ отсутствует), профиль \eqref{eq:profile} ($\rho$ --- стандартное
представление $O(m)$), локальный квадратичный джет
($\rho=\mathrm{std}\oplus\mathrm{Sym}^2$), поле подпространств.

\subsection{Модификации многообразия: карты и $\varphi$-семейство}\label{sec:retr}

\begin{definition}[модификация]\label{def:chart}
Модификация $\mathcal M$ в точке $Y$ --- пара (простое линейное пространство
 $E$, отображение $\chi_Y\colon E\to\mathcal M$), $\chi_Y(0)=Y$,
 $D\chi_Y(0)=\mathrm{id}$. Канонические типы:
(а) аффинная карта в амбиент со штрафом ($B=P+\Delta$);
(б) QR/полярная карта $\qf(Y+Z)$;
(в) $\varphi$-деформация экспоненты.
\end{definition}

\begin{proposition}[$\varphi$-семейство ретракций]\label{prop:retr}
Пусть $Z=V\Sigma A\T$ --- тонкое SVD горизонтального $Z$ ($Y\T Z=0$),
и $R^\varphi_Y(Z)=\Exp_Y\bigl(V\varphi(\Sigma)A\T\bigr)$, где $\varphi$ применяется покоординатно. Тогда:
(i) при $\varphi(0)=0$, $\varphi'(0)=1$ ($\varphi$ нечётна, гладка)
 $R^\varphi$ --- ретракция первого порядка;
(ii) она второго порядка тогда и только тогда, когда $\varphi''(0)=0$;
(iii) все ретракции второго порядка дают один и тот же шаг Ньютона;
(iv) $\Exp_Y(Z)=Y+Z-\tfrac12YZ\T Z+O(\|Z\|\F^3)$.
\end{proposition}
\begin{proof}
Из (iv): $R_Y(tZ)=Y+V\varphi(t\Sigma)A\T-\tfrac12YA\varphi(t\Sigma)^2A\top+O(t^3)$.
Тогда $\tfrac{d^2}{dt^2}\big|_0R_Y(tZ)=\varphi''(0)\,V\Sigma^2A\T-Y\Sigma^2A\top$,
а у геодезической второй член равен $-YZ\T Z=-YA\Sigma^2A\T$; откуда (i)--(ii).
(iii) --- стандартно \cite{ams,absilmalick}: направление Ньютона определяется
внутренними градиентом и гессианом.
\end{proof}

Примеры: $\varphi(t)=t$ --- $\Exp$; $\varphi(t)=\arctan t$ --- полярная
 $\qf(Y+Z)$ (углы $\arctan\sigma_\ell$); $\varphi(t)=2\arctan(t/2)$ --- Кэли
\cite{wenyin}. При точном поиске по скаляру выбор $\varphi$ меняет только
параметризацию: множества достижимости совпадают с точностью до границы
карт с ограниченным углом (у полярной $|\theta|<\pi/2$, граница ---
предел $|\eta|\to\infty$).

\begin{proposition}[полярная карта декалирует плоскости]\label{prop:polardecomp}
Для $Z=V\diag(\eta)A\T$, $A\T A=V\T V=\Id_q$, $Y\T V=0$:
\begin{equation}\label{eq:polarplane}
 R_Y(Z)=Y(\Id-AA\T)+YA\,\diag\!\bigl((1+\eta_\ell^2)^{-1/2}\bigr)A\T
 +V\,\diag\!\bigl(\eta_\ell(1+\eta_\ell^2)^{-1/2}\bigr)A\T.
\end{equation}
Каждая плоскость $(Ya_\ell,v_\ell)$ движется независимо; столбцы $Yb$,
 $b\perp\spanop(A)$, не меняются; угол поворота $\theta_\ell=\arctan\eta_\ell$.
\end{proposition}
\begin{proof}
 $M=Y+V\diag(\eta)A\T$ даёт $M\T M=\Id+A\diag(\eta)^2A\T$ (так как $Y\T V=0$),
откуда $(M\T M)^{-1/2}=\Id+A[\diag((1+\eta_\ell^2)^{-1/2})-\Id]A\T$; подстановка.
\end{proof}

\subsection{Плоское разложение шага}\label{sec:cartan}

\begin{definition}[простое вращение]\label{def:plane}
Для $y=Ya$, $\|a\|=1$, $v\perp\spanop(Y)$, $\|v\|=1$:
 $\rot_t(y,v)=\Id+(\cos t-1)(yy\T+vv\T)+\sin t\,(vy\T-yv\T)$;
 $Y(t)=\rot_t(y,v)\,Y=Y(\Id-aa\T)+(Ya\cos t+v\sin t)a\T$,
 $\rank(Y(t)-Y)\le1$.
\end{definition}

\begin{theorem}[плоское разложение; Картан/CS]\label{thm:cartan}
Пусть $0<m<d$, $Z\in T_Y\Gr$ с тонким SVD $Z=V\Sigma A\T$, $q=\rank Z$.
Тогда:
(i) $\Exp_Y(Z)=\prod_{\ell=1}^q\rot_{\sigma_\ell}(Ya_\ell,v_\ell)\,Y$,
причём вращения попарно коммутируют (плоскости ортогональны);
(ii) то же верно для любой $R^\varphi_Y(Z)$ с углами $\varphi(\sigma_\ell)$ --- та же плоскость, изменённая экспонента;
(iii) множество $\{\Exp_Y(V\Theta A\T)\colon\Theta=\diag(\theta_1,\ldots,\theta_q)\}$ --- вполне геодезический тор, период $\pi$ в координатах подпространства;
максимальный тор имеет размерность $\min(m,d-m)$;
(iv) $\rank(\Exp_Y(Z)-Y)\le q$.
\end{theorem}
\begin{proof}
(i)--(ii): кососимметрический генератор
 $\Omega=\sum_\ell\sigma_\ell(v_\ell y_\ell\T-y_\ell v_\ell\T)$ действует в
попарно ортогональных плоскостях, $\Omega Y=Z$, поэтому
 $e^\Omega Y=\prod_\ell e^{\sigma_\ell\Omega_\ell}Y$ и $e^{\sigma_\ell\Omega_\ell}
=\rot_{\sigma_\ell}(y_\ell,v_\ell)$; подробно \cite[\S2.5]{eas}.
(iii) --- стандартная структура максимальных плоских подмногообразий
компактных симметрических пространств \cite{helgason,gvl}; (iv) ---
 $\Exp_Y(Z)-Y=\bigl(YA(\cos\Theta-\Id)+V\sin\Theta\bigr)A\T$.
\end{proof}

Таким образом, «разложение изменённого многообразия на простое» ---
это разложение Картана: шаг = тор, координаты тора = углы, а всякая
 $\varphi$-модификация экспоненты действует только на параметризацию углов.

\begin{remark}
Жадная траектория (пересчёт градиента после каждого поворота) не совпадает
с одной геодезической: это метод якобиевского типа, а не движение по
заранее выбранному тору.
\end{remark}

\subsection{Спектральный бюджет: какие плоскости выбирать}\label{sec:budget}

\begin{proposition}[наискорейшие направления]\label{prop:budget}
Пусть $G=\grad\Phi(Y)$, $G=U\Sigma V\T$ --- компактное SVD,
 $q'=\rank G\le\min(m,d-m)$.
\begin{enumerate}
\item \textbf{Ранговая шапка + норма Фробениуса:}
 $\min_{\rank Z\le q,\,\|Z\|\F\le1}\ip{G}{Z}=-\|[G]_q\|\F$ достигается на
 $Z_*=-[G]_q/\|[G]_q\|\F$ (фон Нейман + Коши--Буняковский).
\item \textbf{Шар Шаттена-$p$:} при $1/p+1/q=1$  $Z_*\propto-U\Sigma^{q-1}V\T$. Шар $S_2$ даёт градиентный шаг; шар
 $S_\infty$ --- полярный фактор $-UV\T$ (ср. Muon\footnote{K. Jordan et al.,
Muon, 2024: полярный фактор момента как направление спуска.});
шар $S_1$ --- ведущий атом $-u_1v_1\T$ (ваш rank-one).
\item \textbf{Метрика $\mathcal L$ (fixed-$g$ квадратика):} полный шаг в шаре
 $\langle Z,\mathcal LZ\rangle\le1$ есть $Z_*=-\mathcal L^{-1}Q_B$ (Ньютон/GN); для \emph{ранг-один} атома
\begin{equation}\label{eq:Lball}
 \Gamma_1(Q_B)=\max_{a,v}\frac{(a\T Q_Bv)^2}
 {\sum_jN_j(a\T g_j)^2\,v\T(U_j\T U_j)v+\lambda\|a\|^2\|v\|^2},
\end{equation}
замкнутой формы нет --- её и ищет попеременный поиск; при surrogate
 $\mathcal L_0$ с $\rho=\|\mathcal L_0^{-1/2}(\mathcal L-\mathcal L_0)
\mathcal L_0^{-1/2}\|_2<1$ атом из $\mathcal L_0$-преобразованной невязки
даёт долю $(1-\theta)/(1+\theta)$ оптимума, $\theta=\rho/\lambda_{\min}$-типа.
\end{enumerate}
\end{proposition}

\begin{remark}[исправление]\label{rem:kyfan}
Утверждение «шар Ky Fan-$k$ даёт rank-$k$ шаг» неверно: линейный функционал
на шаре Ky Fan концентрируется на ведущем сингулярном направлении
(получается rank-one). Rank-$q$ шаг возникает из \emph{ранговой шапки} с
фробениусовой нормировкой --- п.~1 выше (спектральный бюджет).
\end{remark}

\begin{proposition}[гарантия спуска]\label{prop:descent}
Пусть вторая производная $\Phi$ вдоль единичных геодезических ограничена
константой $L>0$ (существует: $\Gr$ компактен, $\Phi$ гладко при $\tau>0$).
Для $Z=-[G]_q/\|[G]_q\|\F$ и $t=\|[G]_q\|\F/L$:
\begin{equation}\label{eq:descent}
 \Phi(Y)-\Phi(R_Y(tZ))\ge\frac{\|[G]_q\|\F^2}{2L}\ \ge\
 \frac{\|G\|\F^2}{2L\,\min(m,d-m)},
\end{equation}
где использовано $\|G\|\F^2=\sum_{\ell\le q'}s_\ell^2\le\min(m,d-m)\,s_1^2$.
\end{proposition}

\subsection{Точные скалярные профили}\label{sec:angles}

\begin{proposition}[значение профиля в плоскости]\label{prop:angleprofile}
Пусть $Y(t)$ из Опр.~\ref{def:plane}, $A_\perp\in\R^{m\times(m-1)}$,
 $[A_\perp,a]\in O(m)$, $C_j=U_jYA_\perp$, $u_j=U_jYa$, $v_j=U_jv$,
 $\Pi_j=\Id-C_j(C_j\T C_j+\tau\Id)^{-1}C_j\T$. Тогда
\begin{equation}\label{eq:angleprofile}
 \min_g\bigl\{\|I_j-U_jY(t)g\|^2+\tau\|g\|^2\bigr\}
 =k_j-\frac{(\alpha_j\cos t+\beta_j\sin t)^2}
 {a_j\cos^2t+2b_j\sin t\cos t+e_j\sin^2t+\tau},
\end{equation}
где $k_j=I_j\T\Pi_jI_j$, $\alpha_j=u_j\T\Pi_jI_j$, $\beta_j=v_j\T\Pi_jI_j$,
 $a_j=u_j\T\Pi_ju_j$, $b_j=u_j\T\Pi_jv_j$, $e_j=v_j\T\Pi_jv_j$.
 $\Pi_j$ не строится: $x\T\Pi y=x\T y-(C_j\T x)\T(C_j\T C_j+\tau)^{-1}(C_j\T y)$.
Штраф $\lambda d_{\rm ch}^2(Y(t),Y_0)$ --- тригонометрический многочлен
степени 2 (подготовка $O(m^2+dm)$). После подготовки одна оценка угла стоит
 $O(J)$, память $O(1)$ на центр.
\end{proposition}
\begin{proof}
Разложение $g=A_\perp g_\perp+a\gamma$ даёт
 $U_jY(t)g=C_jg_\perp+q(t)\gamma$, $q(t)=u_j\cos t+v_j\sin t$.
Минимизация по $g_\perp$ --- дополнение Шура (значение
 $b\T\Pi_jb$), затем по $\gamma$ --- скалярный минимум квадратичной формы.
\end{proof}

\begin{proposition}[рациональный (полярный) профиль]\label{prop:rational}
В полярной карте $Y(\eta)=\qf(Y+\eta va\T)$ те же шесть скаляров дают
\begin{equation}\label{eq:anglepolar}
 \phi_j(\eta)=k_j-\frac{(\alpha_j+\beta_j\eta)^2}
 {(e_j+\tau)\eta^2+2b_j\eta+(a_j+\tau)},\qquad \eta\in\R,
\end{equation}
без тригонометрии и периодичности; критические точки одного центра ---
корни кубического уравнения. (В тригонометрической карте --- многочлен
степени $\le7$ при замене $z=\tan(t/2)$.)
\end{proposition}

\begin{remark}
\eqref{eq:angleprofile} и \eqref{eq:anglepolar} --- замкнутые формы
\emph{значения}, а не аргминимума: сумма по $j$ рациональных функций может
иметь много минимумов, универсальной замкнутой формы аргминимума нет.
Практика: сетка + Brent (эвристика глобального поиска) либо плоский
trust-region (§2.6). Период в координатах подпространства равен $\pi$.
\end{remark}

\begin{proposition}[плоский refit при $q>1$]\label{prop:flatfit}
При кэше $U_j[Y,V]$ ($m+q$ столбцов) значение
 $\Phi\bigl(\Exp_Y(V\Theta A\T)\bigr)$ для любого набора углов вычисляется
точно малой МНК: $U_jY(\theta)g=[U_jY,\;U_jV]\,\kappa_j(\theta,g)$ с
 $\kappa_j=\bigl[(\Id-AA\top+A\cos\Theta A\top)g;\ \sin\Theta\,A\top g\bigr]$;
стоимость $O\bigl(J(p(m+q)^2+(m+q)^3)\bigr)$ на оценку. Якобиан по углам
также берётся из кэша:
 $\partial_{\theta_\ell}Y(\theta)=(-\sin\theta_\ell\,Ya_\ell+\cos\theta_\ell\,v_\ell)a_\ell\T$.
Производные локальных коэффициентов:
\begin{equation}\label{eq:jacobiandg}
 A_j\dot g_j=T_j\T r_j-M_j\T T_jg_j,\qquad
 \dot r_j=-T_jg_j-M_j\dot g_j,\qquad T_j=U_jZ;
\end{equation}
отбрасывание $\dot g_j$ даёт аппроксимацию Кауфмана --- допустимо только
под контролем убывания.
\end{proposition}

\subsection{Плоский trust-region и сходимость}\label{sec:conv}

Плоский TR: минимизация точного профиля (данные из кэша + точный
тригонометрический штраф) по $\theta\in\mathbb T^q$ методом GN/TR
с принятием шага по \emph{истинному} $\Phi$. Стандартные результаты
о TR на многообразиях \cite{abg} дают глобальную сходимость к
стационарным точкам.

\begin{proposition}[сходимость]\label{prop:conv}
Пусть $\tau>0$, статистики фиксированы; шаги выбираются по
Предл.~\ref{prop:budget}(1) с длиной из \eqref{eq:descent} либо точным
1D-поиском, либо Armijo. Тогда
 $\Phi(Y_{k+1})\le\Phi(Y_k)-c\|G_k\|\F^2$,
 $\sum_k\|G_k\|\F^2<\infty$, все предельные точки стационарны;
по неравенству Лоясевича \cite{ama05} вся последовательность сходится,
а длина траектории конечна.
\end{proposition}

\begin{remark}
При изменении $h,D_j,S_j$ меняется сама цель; гарантии относятся к
внутренней задаче с замороженными статистиками. Внешний цикл
(структурная адаптация $\TEDR_k$) --- неподвижная точка без гарантии
монотонности по $\Phi$.
\end{remark}

\subsection{Алгоритм ППС}\label{sec:alg}

\begin{center}
\fbox{\begin{minipage}{0.96\linewidth}
\textbf{Профильно-плоский шаг} (вход: $Y$, $\{I_j,U_j,N_j\}$, $\tau>0$,
 $\lambda$, $\varepsilon$, $q_{\max}$).
\begin{enumerate}
\item Профиль: $g_j^\star$, $r_j$; кэш $C_j=U_jY$ --- проход $O(Jpdm)$.
\item Градиент $G$ (действия, без $d\times d$) --- проход $O(Jpd)$.
\item Тонкое SVD $G$ ($O(dm^2)$); выбрать $q$ по энергобюджету
 $\sum_{\ell\le q}s_\ell^2\ge(1-\varepsilon)\|G\|\F^2$, $q\le q_{\max}$.
\item Кэш $U_jV$ --- проход $O(Jpdq)$, память $Jpq$.
\item Плоский TR/GN по $\theta\in\mathbb T^q$ с точным профилем
(Предл.~\ref{prop:flatfit}); допуск итераций --- из условия неточного
спуска \eqref{eq:inexact}.
\item Ретракция $R^\varphi_Y$ (по умолчанию полярная:
декалирование, Предл.~\ref{prop:polardecomp}); принятие по истинному
 $\Phi$; при отказе --- деление шага.
\item Численная стабилизация: периодический тонкий QR (не реконструкция
направлений).
\end{enumerate}
Итого на итерацию: 2 прохода по данным, память сверх входа $O(Jp(m+q))$.
\end{minipage}}
\end{center}

\section{Частные случаи}\label{sec:instances}

\begin{center}\small
\begin{tabularx}{\linewidth}{@{}p{30mm}LL@{}}
\toprule
Элемент схемы & SVD-решатель (fixed-$g$) & Профильный Gr-метод\\
\midrule
Многообразие & $\R^{m\times d}$ (плоское) & $\Gr(m,d)$\\
Калибровка & нет ($GL$-инвариантность нарушена штрафом) & $O(m)$-фактор,
профиль по $g_j$\\
Модификация & аффинная карта $B=P+\Delta$, ridge & ретракция $R^\varphi$,
штраф $d_{\rm ch}^2$\\
Атомы & $E=av\T$ & плоскости $(Ya,v)$\\
Выбор атома & SVD $Q_B$ / ALS по $\Gamma_1$ & SVD $G$ (или GN-метрика)\\
Скалярный поиск & $t_*=R/D$ (точно) & профиль \eqref{eq:angleprofile}/%
\eqref{eq:anglepolar}\\
Совместный refit & $M\sigma=h$ в оболочке атомов & $(m+q)$-колоночный
профиль на $U_j[Y,V]$\\
Точное решение & Эккарт--Янг при $\mathcal LX=GXH$ & полный шаг
 $q=\min(m,d-m)$\\
Сертификат & $\|Q_B\|\F^2/\lambda$ & убывание $\Phi$, TR\\
\bottomrule
\end{tabularx}
\end{center}

\subsection{SVD-решатель как евклидов частный случай}\label{sec:svd}

Fixed-$g$ задача: $F(B)=\sum_jN_j\|I_j-U_jB\T g_j\|^2+\lambda\|(B-P)H^{1/2}\|\F^2$,
\begin{equation}\label{eq:Lop}
 \mathcal LX=\sum_jN_j(g_jg_j\T)X\,H_j+\lambda X,\qquad
 C=\sum_jN_jg_jI_j\T U_j+\lambda P,\qquad
 F=C_0+\ip{B}{\mathcal LB}-2\ip{C}{B}.
\end{equation}
Это \eqref{eq:master} с тривиальной калибровкой и пустым $c$: ППС на
 $(\R^{m\times d},\|\cdot\|\F)$. Все перекрёстные блоки по $g_j$ сохраняются; решение строк независимо допустимо лишь при дополнительной
разделимости. При SPD-$H$ замена
 $\widetilde B=BH^{1/2}$, $\widetilde U_j=U_jH^{-1/2}$ точна и сохраняет
ранговые ограничения.

\begin{proposition}[геометрическое вложение]\label{prop:embedding}
Пусть $\Delta=\eta\,va\T$ горизонтальна, $Y_+=\qf(Y+\Delta)$. Тогда
 $Y_+=Y(\Id-aa\T)+\tfrac{Ya+\eta v}{\sqrt{1+\eta^2}}a\T$ ---
та же окружность в $\Gr$, что и rank-one геодезический шаг
\eqref{eq:rotone} с углом $\arctan\eta$; $\rank(Y_+-Y)\le1$.
Кроме того, для горизонтального $\Delta$:
 $d_{\rm ch}^2(\qf(Y+\Delta),Y)=\|\Delta\|\F^2+O(\|\Delta\|\F^4)$,
то есть амбиентный штраф Фробениуса --- модификация chordal-штрафа,
совпадающая с ним до второго порядка.
\end{proposition}
\begin{proof}
Первое --- прямое вычисление из $\Id+\eta^2aa\T$-нормировки
(эквивалент \eqref{eq:polarplane} при $q=1$). Второе:
углы полярной карты $\theta_\ell=\arctan\sigma_\ell$,
 $d_{\rm ch}^2=\sum_\ell\sin^2\theta_\ell=\sum_\ell\sigma_\ell^2/(1+\sigma_\ell^2)$.
\end{proof}

\begin{corollary}
Corrrection-режим SVD-солвера (rank-one $\Delta$ + ортонормировка) и
rank-one Grassmann-поворот движутся по одной и той же окружности и
отличаются только параметризацией угла; все различия методов сосредоточены
в (а) fixed-$g$ против профилированной цели, (б) выборе штрафов,
(в) выборе плоскостей.
\end{corollary}

\begin{theorem}[изотропный одноточечный случай]\label{thm:iso}
Если $\mathcal LX=GXH$, $G,H\succ0$ (в частности
 $H_j=\eta_j\Id$, $G=\sum_jN_j\eta_jg_jg_j\T+\lambda\Id_m$), то
\begin{equation}\label{eq:isoopt}
 B_r=G^{-1/2}\bigl[G^{-1/2}CH^{-1/2}\bigr]_rH^{-1/2}
\end{equation}
глобально минимизирует $F$ при $\rank B\le r$; зазор равен сумме квадратов
отброшенных сингулярных значений $G^{-1/2}CH^{-1/2}$. Для
 $\rank(B-P)\le r$ то же с $C_\Delta=C-\mathcal LP$.
\end{theorem}
\begin{proof}
Замена $Z=G^{1/2}BH^{1/2}$ даёт $\|Z-W\|\F^2$ с
 $W=G^{-1/2}CH^{-1/2}$; обратимые множители сохраняют ранг; теорема
Эккарта--Янга \cite{eckart}.
\end{proof}

Здесь «весь функционал диагонализуется в один flat» --- плоское разложение
выполняется сразу для всей задачи. Сертификаты полного условного решения:
 $\|B-B_*\|\F\le\|Q_B\|\F/\lambda$,
 $F(B)-F(B_*)\le\|Q_B\|\F^2/\lambda$. Условие неточного спуска для
приближённой поправки $D$ с $s=\mathcal LD-Q_B$:
\begin{equation}\label{eq:inexact}
 F(B+D)-F(B)=-\ip{D}{\mathcal LD}+2\ip{D}{s}
 \le-\lambda\|D\|\F^2+2\|s\|\F\|D\|\F,
\end{equation}
достаточное условие убывания: $2|\ip{D}{s}|<\ip{D}{\mathcal LD}$ ---
согласование допусков Крылова с реальным шагом. Ограничение ранга самой
 $B$ принципиально: $\min_{\rank B\le r}\|B-P\|\F^2=m-r$, т.е.
 $F(B)\ge\lambda(m-r)$; недостающие направления из одной $B_r$ не
восстанавливаются --- аргумент в пользу профилированной постановки и
correction-режима. Для $\|\Delta\|_2<1$:
 $d_{\rm ch}\le\|\Delta\|\F/(1-\|\Delta\|_2)$.

\subsection{Жадный rank-one как якобиев метод}\label{sec:rankone}

Последовательность rank-one поворотов (адаптивный выбор плоскости по
текущему градиенту, точный угол) --- покоординатный спуск по торам,
прямой родственник якобиевых методов \cite{jacobi,schonhage} и GROUSE
\cite{grouse}; в амбиенте --- rank-one matrix pursuit \cite{pursuit,mallat}.
Гарантии: монотонность и сертификаты --- §2.4, §3.1; для profiled-цели ---
Предл.~\ref{prop:descent}, \ref{prop:conv}. Слабости: потеря
 $\min(m,d-m)$; при $m=1$ совпадает с градиентным спуском по геодезической.
Оправдан: огромный/разреженный $d$ ($v$ из top-$k$ координат $Ga$),
стриминг, концентрация $s_1^2/\|G\|\F^2$; число шагов тогда служит
неявной регуляризацией (бустинг на $\Gr$). В основном пути ---
абляция.

\subsection{Полный шаг и семейство ретракций}\label{sec:fullstep}

 $q=\min(m,d-m)$: полный геодезический шаг замкнутой формы
\begin{equation}\label{eq:rotone}
 \Exp_Y(Z)=Y(\Id-AA\T)+YA\cos\Theta\,A\T+V\sin\Theta\,A\T,
\end{equation}
углы --- главные углы; совместная оптимизация углов --- Предл.~\ref{prop:flatfit}.
Выбор $\varphi$ (Exp / полярная / Кэли) не меняет оцениватель при точном
поиске по скаляру и не меняет шаг Ньютона; полярная карта дешевле
(без тригонометрии) и декалирует плоскости (Предл.~\ref{prop:polardecomp}).

\subsection{Анизотропный тензор и $\kappa$-семейство метрик}\label{sec:tensor}

Пара $(Y,\Lambda)$ с калибровкой $(YQ,\,Q\T\Lambda Q)$ задаёт объект
 $Y\Lambda Y\T$ --- PSD-матрицу ранга $m$; локализация
 $D=\alpha^2\Id+Y\Lambda Y\T$ (или $D_{\rm orth}=\alpha^2(\Id-YY\T)+Y\Lambda Y\T$)
применяется действиями: $f(D)x=f(\alpha^2)x+Y\bigl(f(\alpha^2\Id+\Lambda)-
f(\alpha^2)\Id\bigr)Y\T x$, стоимость $O(dm)$. Однопараметрическое
семейство целей:
\begin{equation}\label{eq:kappa}
 \Phi_\kappa(Y)=\sum_jN_j\min_g\,(\hat\nabla_j-Yg)\T\hat\Sigma_j^\kappa
 (\hat\nabla_j-Yg)+\lambda d_{\rm ch}^2.
\end{equation}
\begin{proposition}[три тождества]\label{prop:kappa}
(i) $\|I_j-U_jYg\|^2=(\hat\nabla_j-Yg)\T\,\Sigma_jS_j\T S_j\Sigma_j\,
(\hat\nabla_j-Yg)$; при $S_j=\Id$ метрика $=\Sigma_j^2$ ($\kappa=2$).
(ii) $\sum_iw_{ij}(y_i-\bar y_j-z_{ij}\T Yg)^2
=\mathrm{const}_j+N_j(\hat\nabla_j-Yg)\T\Sigma_j(\hat\nabla_j-Yg)$ ---
локальный критерий MAVE \cite{mave} ($\kappa=1$) при тех же центрах,
весах и intercept.
(iii) При $\kappa=0$, $\tau=0$ профиль равен
 $\mathrm{const}-\tr(Y\T C_{\rm grad}Y)$,
 $C_{\rm grad}=\sum_jN_j\hat\nabla_j\hat\nabla_j\T$ --- точный глобальный
OPG-оптимум: ведущие $m$ собственных векторов.
\end{proposition}
При независимом шуме $\operatorname{Cov}(c_j\,|\,x,w)
=\frac{\sigma^2}{N_j^2}\sum_iw_{ij}^2z_{ij}z_{ij}\T$; для $w\in\{0,1\}$ это
 $\sigma^2\Sigma_j/N_j$ --- GLS-метрика $\Sigma_j$, т.е. $\kappa=1$ (эффективный выбор). Идентифицируемость равносильна $H_{GN}\succ0$ (градиенты натягивают $m$-мерное подпространство); в истинной точке
GN-гессиан равен информации Фишера, асимптотическая ковариация в
касательном пространстве $\approx\sigma^2H_{GN}^{-1}$ --- при корректной
гауссовой модели; к исходной finite-sketch цели это без условий не
относится. Конус PSD ранга $m$ --- симметрическое пространство: его flat
--- диагональные $\Lambda$, простые движения --- собственные плоскости;
та же ППС применима к внешнему циклу (без гарантии монотонности
неподвижной точки по $h$).

\subsection{Поле подпространств}\label{sec:field}

Замена $Y$ на $Y_j$ без связи даёт $J$ независимых задач (гессиан
блочно-диагонален). Геометрически корректное поле малоранговых шагов ---
Tucker-параметризация \emph{касательного поля} до ретракции:
\begin{equation}\label{eq:tangentfield}
 Z_j=V\Bigl(\sum_{\ell=1}^sw_\ell(x_j)M_\ell\Bigr)A\T,\qquad
 Y_j=R_{Y_0}(Z_j),
\end{equation}
общая оболочка $\spanop[Y_0,V]$ размерности $m+q>R$-ограничений;
хранение $O(d(m+q)+mk+sqk+Js)$ вместо $O(Jdm)$. Пространственная
гладкость --- через $w_\ell(x)$ (графовый штраф $\tr(w\T Lw)$), штраф
соседства --- инвариантно: $\sum\omega_{ij}d_{\rm ch}^2(Y_i,Y_j)$.
CP/Tucker в терминах самих $B_j$ --- та же идея с оговоркой: при
 $\rank B_j=m=R$ поле не вращается; нужно $R>m$. Интегрируемость поля
плоскостей (инволютивность) --- отдельное условие, не требующееся для
регрессии.

\section{Научные предшественники}\label{sec:lit}

\begin{center}\small
\begin{tabularx}{\linewidth}{@{}p{34mm}L@{}}
\toprule
Направление & Что общего; где заканчивается перенос\\
\midrule
Якобиевы методы \cite{jacobi,schonhage,bgbm} &
Жадные/циклические плоские вращения, квадратичная сходимость циклической
схемы. У нас --- профилированная нелинейная цель и точный угловой поиск,
а не внедиагональная норма.\\
CS-разложение, главные углы \cite{gvl} &
Плоские координаты $\Gr$ --- тор главных углов; основа Теоремы~\ref{thm:cartan}.\\
Симметрические пространства, Картан \cite{helgason,eas} &
 $\mathfrak g=\mathfrak k\oplus\mathfrak p$, максимальный flat, вполне
геодезические торы; атомы --- корневые плоскости.\\
GROUSE \cite{grouse}; Oja \cite{oja} &
Rank-one геодезические обновления, стриминг. У нас --- профиль, точный
линейный поиск, произвольные операторы $U_j$, а не выборка координат.\\
VARPRO \cite{varpro,varpro2} &
Профилирование и производные; наши замкнутые скалярные профили ---
VARPRO с одним нелинейным параметром на центр.\\
Ретракции \cite{absilmalick,wenyin,ams} &
Семейство $\varphi$, порядок ретракций, единый Ньютон.\\
TR на многообразиях \cite{abg}; Лоясевич \cite{ama05} &
Глобальная сходимость к стационарным точкам; сходимость траектории.\\
Наискорейший спуск в нормах Шаттена; Muon (2024) &
Полярный фактор как направление; у нас --- частный случай спектрального
бюджета (Предл.~\ref{prop:budget}(2)).\\
Rank-one pursuit, OMP \cite{pursuit,mallat} &
Жадные атомы в амбиенте с точным масштабом; гарантии matrix completion
на ADP не переносятся.\\
Эккарт--Янг \cite{eckart}; weighted low-rank \cite{sw,mishra} &
Замкнутая форма при $GXH$; в общем случае взвешенная малоранговая
аппроксимация требует иной оптимизации.\\
Координатный спуск \cite{tseng}; дефляция/ coarse space \cite{defcg} &
Циклическая минимизация по углам тора; flat как coarse space.\\
MAVE/OPG/SIR/HJPS \cite{mave,sir,hjps} &
Статистические цели; $\kappa$-семейство связывает их алгебраически
(Предл.~\ref{prop:kappa}).\\
Envelope-модели \cite{cook} &
Глобальная низкоразмерная оболочка; аналог --- envelope $V$ в поле.\\
Случайное SVD \cite{halko} &
Только при дорогих действиях; при малом $m$ тонкое SVD $d\times m$ дёшево.\\
\bottomrule
\end{tabularx}
\end{center}
Целостная комбинация «профиль + плоское разложение + точный угловой поиск +
спектральный бюджет» в известной нам литературе как единый метод не
представлена; ближайшие точечные прототипы --- GROUSE (без профиля и
точного поиска) и якобиевы методы (без фактор-многообразия). Новизна
относительно всей литературы и преимущество в восстановлении здесь не
доказаны.

\section{Улучшения}\label{sec:improve}

\textbf{И1. Плоский шаг ранга $q$ вместо жадного rank-one} (основной путь).
Один-два прохода по $U_j$ на итерацию против $\sim$40 проходов на атом
жадной схемы (ALS с matvec $O(Jpd)$); снимается потеря
 $\min(m,d-m)$ из \eqref{eq:descent}; убывание $\ge\|[G]_q\|\F^2/(2L)$;
сходимость --- Предл.~\ref{prop:conv}. [формулы доказаны; wall-clock и
качество --- замерить]

\textbf{И2. Предобусловленный выбор плоскостей.} Метрика шага
 $\mathcal M_Y$ = GN-гессиан (или surrogate $\mathcal L_0$ с
 $\rho$-границей); плоскости --- по $\mathcal M_Y^{-1}$-преобразованному
градиенту (натуральный градиент/GN); сертификат доли выигрыша через
 $\rho$. Эффект --- качество при сильной анизотропии после структурной
адаптации. [доказаны оценки; измерить]

\textbf{И3. Полярно-рациональный профиль.} $\varphi=\arctan$: на центр ---
кубика \eqref{eq:anglepolar}, без тригонометрии и периодичности; та же
память (шесть скаляров). Эффект --- скорость и устойчивость углового
поиска. [доказано; замерить]

\textbf{И4. Статистический контур: $\kappa$-гомотопия.}
Инициализация $\kappa=0$ (OPG, замкнуто, Предл.~\ref{prop:kappa}(iii));
далее $\kappa\colon0\to1$ (MAVE/GLS-эффективная метрика) вместе с
 $h_k\downarrow$; детерминированный GLS-выбор $S_j=\hat\Sigma_j^{-1/2}$;
chordal-штраф; $\tau>0$ для аналитичности. Эффект --- качество
восстановления. [тождества доказаны; estimator-часть --- измерить]

\textbf{И5. Двухуровневая доводка.} После плоского шага --- CG/TR в
касательном пространстве с дефляцией плоских направлений \cite{defcg};
допуск Крылова --- из \eqref{eq:inexact}; принятие по истинному $\Phi$ гарантирует неувеличение. Снимает внеплоскостную кривизну малым числом
matvec'ов. [предложение]

\textbf{И6. Память.} Тонкое QR-сжатие $[U_j,I_j]=Q_jR_j\to(\bar U_j,\bar I_j)$ без усечения --- точное тождество (выигрыш при $p>d+1$); факторизованные/
потоковые $U_j$; кэш $Jp(m+q)$; батчинг $U_j[W]$; запрет плотных $d\times d$.
[точно/численно]

\textbf{И7. Тензор: ППС на конусе PSD.} Диагональ $\Lambda$ --- flat,
собственные плоскости --- атомы; та же схема для внешнего цикла; обновление
 $(Y,\Lambda)$ с согласованной калибровкой. [предложение]

\textbf{И8. Поле: общие плоскости + гладкие веса.} \eqref{eq:tangentfield}
с $R=m+q$, графовой гладкостью $w_\ell$ и инвариантным штрафом соседства;
идентифицируемость поля --- отдельно. [предложение]

\medskip
\textbf{Порядок внедрения.} (1) OPG-инициализация (И4, $\kappa=0$);
(2) ППС с адаптивным $q$ и TR по углам (И1, И3); (3) гомотопия $\kappa$ (И4); (4) предобусловливание и двухуровневая доводка (И2, И5); (5) память
(И6); (6) тензор/поле (И7, И8). Baseline сравнения: текущий AO, fixed-$g$ SVD, MAVE, OPG, SIR; метрика $d_{\rm ch}^2(Y,Y_*)$, wall-clock, память;
протокол --- из \texttt{EDR\_unified\_theory} (paired frozen statistics,
далее полные фиты по заранее заданным seeds).

\section{Границы результата}\label{sec:limits}

Не доказано и не заявляется: глобальная оптимальность плоских шагов и
жадных схем (кроме изотропного случая, Теорема~\ref{thm:iso}); скорость
сходимости (только асимптотика); монотонность внешнего цикла адаптации
ядра; оптимальность адаптивного $q$; глобальность сетки+Brent
(эвристика; сертификат требует границ на интервале либо перечисления
критических точек); идентифицируемость поля. Все новые градиенты и
Якобианы проверять конечными разностями; численные проверки углового
профиля и ретракций из исходных записей (совпадение до $10^{-15}$)
согласуются с Предл.~\ref{prop:angleprofile}, \ref{prop:rational}.

\appendix
\section{Карта исходных записей и исправления}
\begin{center}\small
\begin{tabularx}{\linewidth}{@{}p{36mm}L@{}}
\toprule
Материал & Место и уточнения\\
\midrule
\texttt{multiindex.tex} &
Модель, статистики, AO, структурная адаптация: §\ref{sec:model},
§\ref{sec:tensor}. Single-index --- $m=1$.\\
\texttt{SVD\_solver.tex}, \texttt{SVD\_algorithm.tex}, \texttt{SVD\_cur.tex} &
Оператор, жадные атомы, изотропная форма, сертификаты, rank floor:
§\ref{sec:svd}. Формализовано как евклидов частный случай ППС;
вложение --- Предл.~\ref{prop:embedding}.\\
\texttt{grassmann\_rank\_one.tex} &
Профиль, градиент, rank-one экспонента: §\ref{sec:cartan}, §\ref{sec:rankone}.
Уточнено: «Taylor-семейство» --- семейство карт $\varphi$, а не разложение
функции; жадная траектория $\ne$ одна геодезическая; добавлена сходимость.\\
\texttt{EDR\_unified\_theory.tex} &
Мастер-класс, четыре метрики, поле, стоимости: §\ref{sec:master}--%
\ref{sec:field}, §\ref{sec:improve}. Добавлены полярная декалировка
(Предл.~\ref{prop:polardecomp}) и рациональный профиль
(Предл.~\ref{prop:rational}).\\
\texttt{Tucker.md} &
Поле и envelope: §\ref{sec:field}. Исправлен знак ridge-штрафа
(положительный); low rank третей моды сам по себе не даёт гладкости.\\
\texttt{cladue.md} &
Schur-поиск угла, семейства норм, GN, гомотопия: §\ref{sec:angles}--%
\ref{sec:tensor}, §\ref{sec:improve}. Исправлено: шар Ky Fan-$k$ даёт
rank-one, а не rank-$k$ (Замечание~\ref{rem:kyfan}); covariance
 $w_{ij}^2$; граница полярной карты; заявления о глобальном поиске и
сходимости даны с условиями.\\
\bottomrule
\end{tabularx}
\end{center}

\begin{thebibliography}{99}\small
\bibitem{jacobi} C.~G.~J. Jacobi. \emph{\"Uber ein leichtes Verfahren, die in der
Theorie der S\"acularst\"orungen vorkommenden Gleichungen numerisch aufzul\"osen}.
J. Reine Angew. Math. 30, 51--94, 1846.
\bibitem{schonhage} A. Sch\"onhage. \emph{Zur Konvergenz des Jacobi-Verfahrens}.
Numer. Math. 3, 374--380, 1961.
\bibitem{bgbm} A. Bunse-Gerstner, R. Byers, V. Mehrmann.
\emph{Numerical methods for simultaneous diagonalization}. Linear Algebra Appl. 191, 49--71, 1993.
\bibitem{gvl} G.~H. Golub, C.~F. Van Loan. \emph{Matrix Computations}, 4th ed.
Johns Hopkins University Press, 2013.
\bibitem{helgason} S. Helgason. \emph{Differential Geometry, Lie Groups, and
Symmetric Spaces}. Academic Press, 1978.
\bibitem{eas} A. Edelman, T.~A. Arias, S.~T. Smith.
\emph{The geometry of algorithms with orthogonality constraints}.
SIAM J. Matrix Anal. Appl. 20(2), 303--353, 1998.
\bibitem{ams} P.-A. Absil, R. Mahony, R. Sepulchre.
\emph{Optimization Algorithms on Matrix Manifolds}. Princeton University Press, 2008.
\bibitem{absilmalick} P.-A. Absil, J. Malick.
\emph{Projection-like retractions on matrix manifolds}. SIAM J. Optim. 22(1), 135--158, 2012.
\bibitem{wenyin} Z. Wen, W. Yin.
\emph{A feasible method for optimization with orthogonality constraints}.
SIAM J. Optim. 23(4), 2462--2488, 2013.
\bibitem{ama05} P.-A. Absil, R. Mahony, B. Andrews.
\emph{Convergence of the iterates of descent methods for analytic cost functions}.
SIAM J. Optim. 16(2), 541--557, 2005.
\bibitem{abg} P.-A. Absil, C.~G. Baker, K.~A. Gallivan.
\emph{Trust-region methods on Riemannian manifolds}. Found. Comput. Math. 7(3), 303--330, 2007.
\bibitem{varpro} G.~H. Golub, V. Pereyra.
\emph{The differentiation of pseudo-inverses and nonlinear least squares problems
whose variables separate}. SIAM J. Numer. Anal. 10(2), 413--432, 1973.
\bibitem{varpro2} G.~H. Golub, V. Pereyra.
\emph{Separable nonlinear least squares: the variable projection method and its
applications}. Inverse Problems 19, R1--R26, 2003.
\bibitem{eckart} C. Eckart, G. Young.
\emph{The approximation of one matrix by another of lower rank}. Psychometrika 1, 211--218, 1936.
\bibitem{grouse} L. Balzano, R. Nowak, B. Recht.
\emph{Online identification and tracking of subspaces from highly incomplete
information}. Allerton, 2010; arXiv:1006.4046.
\bibitem{oja} E. Oja. \emph{Simplified neuron model as a principal component analyzer}.
J. Math. Biol. 15, 267--273, 1982.
\bibitem{pursuit} Z. Wang, M.-J. Lai, Z. Lu, W. Fan, H. Davulcu, J. Ye.
\emph{Rank-one matrix pursuit for matrix completion}. ICML (PMLR 32), 91--99, 2014.
\bibitem{mallat} S. Mallat, Z. Zhang.
\emph{Matching pursuits with time-frequency dictionaries}.
IEEE Trans. Signal Process
